\section{Introduction}
\label{sec:intro}

Scaling test-time compute has become the dominant lever for improving large models on artifact-producing tasks: code, web pages, slides, figures, animations, video edits, research reports. The literature has formalized two axes. \emph{Reasoning scaling} \citep{wei2022chain,yao2023tree,deepseek2025r1,snell2024scaling} spends more tokens thinking before committing to an answer. \emph{Sampling scaling} \citep{wang2023selfconsistency,brown2024large} draws more independent attempts and selects one. The two look different, but they share a property that this paper treats as fundamental: both are \textbf{internal}. Every additional token, whether spent on a longer chain of thought or another sample, is generated from the same frozen weights and the same fixed prompt. More internal compute reorganizes information the system already has; it imports none.

\begin{figure}[t]
\centering
\resizebox{\linewidth}{!}{%
\begin{tikzpicture}[font=\small,
  box/.style={rounded corners=2pt,draw=cMut,thick,minimum height=8mm,align=center,fill=white,inner sep=4pt},
  prop/.style={box,fill=cStu!8,draw=cStu!70},
  env/.style={box,fill=black!4},
  gnd/.style={box,fill=cGnd!10,draw=cGnd},
  ung/.style={box,fill=cUng!8,draw=cUng},
  ar/.style={-{Latex[length=2mm]},thick,cMut},
  gar/.style={-{Latex[length=2mm]},very thick,cGnd},
  uar/.style={-{Latex[length=2mm]},thick,cUng,dashed}]
  \node[prop] (prop) {Proposer\\\footnotesize(frozen LLM)};
  \node[env, right=14mm of prop] (art) {Artifact\\\footnotesize code / HTML / SVG};
  % GROUNDED lane (top)
  \node[gnd, right=15mm of art, yshift=9mm] (inst)
     {Instrument\\\footnotesize\emph{execute / render / measure}};
  \node[gnd, right=13mm of inst] (score) {\footnotesize\textbf{(2)} grounded\\\footnotesize evaluation};
  \node[prop, below=7mm of inst, xshift=6mm] (rev) {Reviewer};
  % UNGROUNDED lane (bottom)
  \node[ung, below=15mm of art, xshift=3mm] (vlm) {VLM reads a \emph{screenshot}};
  \node[ung, right=10mm of vlm] (uscore) {\footnotesize blind feedback\\\footnotesize \& blind score};
  % loop arrows
  \draw[ar] (prop) -- node[above,font=\footnotesize]{generate} (art);
  \draw[ar] (art.east) -- (inst.west);
  \draw[gar] (inst) -- (score);
  \draw[gar] (inst.south) -- node[left,font=\footnotesize,cGnd,align=center,pos=0.45]{\textbf{(1)}\\defects} (rev.north);
  \draw[ar] (rev.south) |- ([yshift=-9mm]art.south) -| node[pos=0.25,below,font=\footnotesize]{targeted revision} (prop);
  % ungrounded branch
  \draw[uar] (art.south) -- (vlm.north);
  \draw[uar] (vlm) -- (uscore);
  \begin{scope}[on background layer]
    \node[draw=cGnd!60,dashed,rounded corners,fill=cGnd!4,fit=(inst)(score)(rev),inner sep=3mm,
      label={[cGnd,font=\footnotesize\bfseries]above:GROUNDED: measure the rendered DOM / run \texttt{pytest}}] {};
    \node[draw=cUng!60,dashed,rounded corners,fill=cUng!4,fit=(vlm)(uscore),inner sep=2.5mm,
      label={[cUng!85!black,font=\footnotesize\bfseries]below:UNGROUNDED: overflow cropped off-frame, rates broken figures ``perfect''}] {};
  \end{scope}
\end{tikzpicture}}
\caption{\textbf{Grounding must hold on both sides of the interaction loop.} A proposer emits an artifact, and a grounded instrument \emph{executes, renders, or measures} it. That single observation does double duty. \textbf{(1)} As grounded \emph{feedback}, it becomes a defect list the reviewer turns into a targeted revision---the channel that lets interaction import information from outside the weights and escape the internal ceiling (\Cref{sec:framework}). \textbf{(2)} As grounded \emph{evaluation}, it is the score that makes the gain measurable. Substituting the default VLM-on-a-screenshot judge (orange lane) breaks \emph{both}: a screenshot drops off-canvas overflow and fine overlap before the model ever sees them, so it neither scores the defect nor---as a reviewer---fixes it.}
\label{fig:concept}
\end{figure}

A third axis breaks this closure: \emph{interaction} with an external instrument that observes the artifact itself---run the tests, render the page and measure the layout, issue the search query. Such an observation is not a function of the weights. It reports how the artifact actually behaves, and it can carry information the model never had. This axis is well established empirically \citep{shinn2023reflexion,yao2023react,gou2024critic,shen2025thinking}, but accounts of \emph{why} and \emph{when} it should beat internal scaling remain informal and, we argue, half-stated. Practitioner reports at industrial scale sharpen the stakes: in a controlled $900$-run deployment study at ByteDance, frontier coding models pass the functional-correctness bar yet miss on delivery quality until a feedback ``harness'' is added---and that report had to invent a multi-axis quality metric before the gap was even visible \citep{hong2026aicoding}. That is a field observation of exactly the two-sided problem we formalize below (detailed in \Cref{sec:related}).

\paragraph{Grounding is the variable, on both sides of the loop.} Our thesis is that one variable governs interaction scaling: \textbf{grounding}---the property that a feedback signal is produced by an instrument observing the artifact's actual form or behavior, rather than by a model rendering an opinion about it. And grounding must hold on \emph{both} sides of the proposer--reviewer loop (\Cref{fig:concept}):
\begin{enumerate}[leftmargin=1.4em,itemsep=1pt,topsep=2pt]
\item \textbf{Grounded feedback} (the signal that drives revision): run the code and read the failing test, render the HTML and measure its bounding boxes, issue the query and read the result. This is the side prior work emphasizes, and a one-paragraph information argument (\Cref{sec:framework}) says it is exactly what lets interaction escape the ceiling that bounds reasoning and sampling.
\item \textbf{Grounded evaluation} (the metric that \emph{measures} the gain): the same observation, used as the score. This side is almost always overlooked, and getting it wrong silently hides the effect.
\end{enumerate}
The second point is not a detail. The default evaluator for visual artifacts is a vision--language model (VLM) reading a screenshot, and a screenshot is an instrument whose channel \emph{drops the defect before the judge ever sees it}: content overflowing the canvas is cropped off-frame, and small overlaps fall below the model's visual acuity. On dense academic figures this judge rates $14$ of $15$ single-shot renders ``perfect'' when a deterministic measurement finds only $3$ actually clean (\Cref{sec:grounded_eval}). Worse, when the same VLM is used as the \emph{reviewer}, its edits make slide layouts worse (\Cref{sec:feedback}). Ungrounded on either side, the third axis either does not fire or cannot be seen.


\Needspace*{6\baselineskip}
\paragraph{Contributions.}
\begin{enumerate}[leftmargin=1.4em,itemsep=1pt,topsep=2pt]
\item \textbf{A grounding framework (\Cref{sec:framework}).} A two-category feedback taxonomy---\emph{ungrounded} model opinion vs.\ \emph{grounded} instrument observation---together with a \emph{coverage} principle: grounded feedback helps exactly as far as the instrument's observational reach. A one-paragraph information argument (formal statement in \Cref{app:ceiling}) explains why internal scaling saturates, and a symmetric observation says an ungrounded \emph{metric} cannot measure the gain. Beyond the documented \emph{biases} of model-as-judge evaluation \citep{zheng2023judging}, we isolate a sharper failure: for artifacts whose quality is a measurable physical property, a screenshot judge is \emph{structurally blind}---not merely noisy---which silently hides a real interaction-scaling effect.
\item \textbf{Interaction keeps scaling where internal compute stops (\Cref{sec:interaction}).} At a matched token budget, reasoning-only and best-of-$N$ saturate---the latter even with an oracle verifier---while every interaction strategy approaches a perfect pass rate, the proposer--reviewer harness doing so at the lowest token cost with zero seed variance.
\item \textbf{The feedback instrument must observe the defect (\Cref{sec:feedback}).} Holding the task suite fixed and swapping only the reviewer's instrument: execution feedback converges far cheaper than critique on behavioral bugs; a linter (grounded, but observing only form) buys nothing on them; and a screenshot-reading VLM reviewer makes slide geometry \emph{worse} while a geometric reviewer repairs it.
\item \textbf{So must the metric (\Cref{sec:grounded_eval}).} The standard VLM judge passes $14$ of $15$ broken figures; a deterministic DOM-geometry instrument instead reveals large, statistically decisive defect reductions under grounded feedback across four visual modalities measured identically---gains the standard metric cannot even see.
\item \textbf{Robustness and internalization (\Cref{sec:robustness,sec:distill}).} The effects replicate across three model families (including the geometry result with Gemini~3.1~Pro as proposer), on a held-out suite, and across a budget-allocation sweep. The harness's behavior also partially distills into an $8$B student at ${\sim}10\times$ lower deployment cost.
\end{enumerate}
